% ============================================================================================
\section{Discussion: Cross-Task Insights}
\label{sec:discussion}
\textbf{Heat, moisture and wind on the same dates.} Using the same eight dates in S1--S3 lets us read the three fields
together.
\begin{itemize}
\item \emph{D1 (8~Jul).} The western heat wave ($\Tmax$ up to 53\,\degC, Fig.~\ref{fig:07}) coincides with VPD of
  4--10\,kPa over the same deserts (Fig.~\ref{fig:16}). Meanwhile Beryl's circulation (Fig.~\ref{fig:21}a) pumps saturated
  Gulf air inland. Texas and Arkansas are hot (33--38\,\degC) but have low VPD and heavy rain (Fig.~\ref{fig:09}a). Heat
  stress there is humid heat, not drought stress.
\item \emph{D2 (24~Jul).} The warm anomaly over Montana and the Dakotas (+10 to +15\,\degC, Fig.~\ref{fig:06}) matches the
  season's largest $\ge$2\,kPa VPD area (34.2\,\%, Table~\ref{tab:evol}). The cool anomaly band from Texas to Tennessee is
  where Fig.~\ref{fig:09} shows rain. Temperature anomalies largely follow cloud and rain.
\item \emph{D3 (5~Aug).} Debby's closed vortex over the Big Bend (strongest vorticity of the four storms) coincides with
  the season's wettest cell (331\,mm on D3; 401\,mm on the previous day) and a sub-1\,kPa VPD footprint over the Southeast.
  The slow storm translated little, so wind, rain and moisture overlap.
\item \emph{D7--D8 (26--27~Sep).} Helene drives the largest rain area (93.5\,mm box mean), the strongest daily winds
  (16.1\,m/s), a $-10$ to $-13$\,\degC{} daytime anomaly under its rain shield (Fig.~\ref{fig:06}), and a moist VPD
  footprint over the whole Southeast. Between D7 and D8 the wind vortex moves to Kentucky and Tennessee (Fig.~\ref{fig:24}h)
  and the heaviest rain moves to the southern Appalachians (Fig.~\ref{fig:09}h). This is the vortex-to-rain chain behind
  the inland flooding.
\item \emph{Season.} Mean $\Tmax$ falls by about 4\,\degC{} from D1--D2 to D6--D8, the land share with VPD $\ge$2\,kPa halves,
  and mean wind speed rises by about 40\,\% (2.8 to 3.9--4.0\,m/s). Summer heat gives way to more energetic early-autumn
  circulation.
\end{itemize}

\textbf{Method-level insights.}
(1) Normalisation and level choice repeat across tasks. Global vs.\ local colour scaling (S1), uniform vs.\ quantile vs.\
domain contour levels (S2) and proportional vs.\ uniform arrows (S3) are the same decision: \emph{comparability across
items vs.\ contrast within an item}.
(2) Aggregation is both necessary and a source of error: block averaging of vectors (S3), Gaussian smoothing of contours
(S2), monthly vs.\ yearly bins (I3) and k-means bundles (I2a). In each case we quantified what was lost (aliasing angles,
isoline length, missing turning points, silhouette).
(3) On the InfoVis side, the same structure appears through four lenses. The pandemic peak and the post-2021 decline show
in the treemap (H3), the horizon plots and the year ridgelines. The visual-vs.\ discussion split shows in the treemap
measures (score vs.\ comments), the parallel-coordinates brushing, the comment ridgelines and the network blocs (image
hosts vs.\ news outlets). Each technique alone gives one view; together they triangulate the finding.

\section{Conclusion}
We built a reproducible pipeline: one command per task, every preprocessing decision logged, and every table generated from
data. With it we produced 42 report figures (55 image files including sub-figures, plus five GIFs) covering colour mapping,
contouring with an exactly validated Marching Squares and asymptotic decider, vector-field sampling and streamline seeding,
five network layouts with quantitative comparison, D3 parallel coordinates with bundling, brushing and reordering, a
zoomable treemap with three remodels and five tilings, and hand-built horizon and ridgeline plots. The SciVis results show
how the 2024 hurricanes and the western heat wave appear in, and are distorted by, different encodings of daily 4\,km
fields. The InfoVis results show how Reddit's largest communities divide into link ecosystems and engagement styles, and
how strongly an all-time-top sample is shaped by survivorship. Future work: hourly data (e.g.\ ERA5) to resolve
hurricane winds, canvas or WebGL rendering for parallel coordinates on all 49k posts, and a multi-decade gridMET baseline for
anomalies.
